\chapter{Testing and Results}
\label{ch:testing}

\todo{acceptance checks and metrics are drafted in
docs/03-planning/testing-and-success.md}

\section{Testing Strategy}

\todo{unit, integration, and scenario levels: what each covers and why}

\section{Unit Tests}

\todo{tests for BOM explosion, inventory netting, capacity feasibility, and
offer scoring, with a coverage summary}

\section{Integration Tests}

\todo{tests for a full negotiation round, including message redelivery to prove
idempotency}

\section{Acceptance Criteria}

\todo{map each functional requirement from Chapter~\ref{ch:problem} to the test
that demonstrates it, as a table. Examiners follow this table closely}

\section{Demonstration Results}

\subsection{Baseline Plan}

\todo{the plan produced for the two competing orders: committed material,
reservations, supplier orders, scheduled slots, total cost, completion time}

\subsection{Rejected Alternatives}

\todo{the offers that lost and the reason each lost, as a table. This is the
evidence for the explainability claim}

\subsection{Plan After Disruption}

\todo{the revised plan, which commitments were released, which were preserved,
and the change in cost and completion time}

\section{Behaviour Without the Language Model}

\todo{repeat the scenario with every model provider stopped, and show the plan
is identical. This demonstrates the constraint that the model never decides
anything}

\section{Discussion}

\todo{what the results show, where the approach is weak, and which numbers
would change at larger scale}
